\noindent\textbf{Resumo:}

\noindent Embora a legislação de transparência exija o mapeamento de procedimentos administrativos, sistemas como o Sistema Eletrônico de Informações (SEI) registram ações isoladas sem consolidar o fluxo real dos processos. O problema desta pesquisa consiste em determinar se é possível extrair, a partir de registros transacionais brutos, um catálogo auditável e confiável dos procedimentos efetivamente executados. Trata-se de uma pesquisa aplicada e quantitativa, conduzida como estudo de caso na Secretaria de Administração do Estado do Piauí (SEAD-PI), analisando 5.243 processos concluídos em 2025. Propõe-se um \emph{pipeline} que modela os eventos como um grafo de conhecimento e agrupa os processos combinando o texto do documento inicial ao histórico de tramitação. Para predizer o caminho típico de cada processo, avaliaram-se seis famílias de modelos (incluindo cadeias de Markov e arquiteturas GRU), testadas em três alfabetos de traço: ação por unidade, roteamento e procedimento. Os resultados mostram ganhos preditivos consistentes contra partições-controle embaralhadas em todas as 252 comparações. Comparado a um modelo global, obteve-se ganho preditivo significativo apenas no alfabeto de procedimento (0,109 com Markov de terceira ordem e 0,131 com GRU). As métricas internas de clusterização não apontaram um agrupamento único ideal. Contudo, a aplicação das técnicas gerou um catálogo validado de 60 tipos de processos, com 75,1\% de cobertura no alfabeto de roteamento, atendendo aos requisitos estritos de anonimização. Conclui-se que o uso de inteligência artificial sobre \emph{logs} administrativos supre a lacuna empírica dos sistemas, viabilizando a governança orientada a dados.

\vspace{0.6\baselineskip}
\noindent\textbf{Palavras-chave:} \textit{mineração de processos; inteligência artificial; redes neurais; otimização; transparência administrativa}.

\vspace{1.2\baselineskip}
\noindent\textbf{Abstract:}

\noindent Access-to-information laws presume the existence of consolidated records for administrative procedures, which often are never created. This research investigates whether an auditable catalogue of executed process types and their typical paths can be derived directly from raw event logs. We conducted an applied, quantitative single-case study at the Secretariat of Administration of the State of Piaui, analyzing 5,243 concluded cases from the Electronic Information System (SEI) in 2025. The methodology builds an event knowledge graph, clusters cases by combining opening-document text with routing and action history, and fits typical paths using six model families across three trace alphabets (action-by-unit, routing, and procedure). Results show positive predictive gains over shuffled control partitions in all 252 baseline comparisons. Compared to a single global model, predictive gains were significant only in the procedure alphabet (0.109 for a third-order Markov chain; 0.131 for GRU); in action-by-unit and routing, gains were indistinguishable from residuals. Although accuracy separated the 14 candidate clusterings by only 0.0055 and internal indices did not reliably predict downstream quality, the pipeline successfully delivered a catalogue of 60 process types with 0.751 coverage in the routing alphabet. Furthermore, the required automated anonymization proved non-inferior at the pre-registered margin across all 54 comparisons.

\vspace{0.6\baselineskip}
\noindent\textbf{Keywords:} \textit{process mining; clustering; graphs; administrative transparency}.

\newpage